\chapter{第十三讲\quad 罪数}
\label{lecture:13}
\noindent{\small\hyperref[ref:lecture:13]{查阅本讲扩展规则}}\par

\section{第一节\quad 一个行为}
\paragraph{先确定评价单位}
令 \(U\) 为依构成要件归并后的行为单位集合，\(E_i(u)\) 表示同一人对同一单位 \(u\) 满足罪名 \(i\) 的构成及相应形态。划分单位时标出实行过程、对象、目的、风险和法条的一体评价要求。自然动作次数、受害人个数和罪名个数分属三种计数。持续看守同一被拘禁者数小时，可构成一个持续的拘禁单位。拘禁行为持续到释放时；已经既遂、尚未终了时，后来加入者的故意参与可按共同犯罪规则另行评价。第89条相应把继续状态的时效起点放在行为终了时。若伤害已经终了，后来仅见到伤者伤势持续，不能据此把已经结束的伤害视为仍在实行。

\paragraph{法条关系的三种测试}
固定行为、对象、解释口径和形态范围，写 \(S_i=\{u:E_i(u)\}\)。区分三个问题：
\[
\begin{aligned}
 S_i\subsetneq S_j &\quad\text{特别要件包含一般要件};\\
 S_i\cap S_j=\varnothing &\quad\text{给定评价单位内两要件互斥};\\
 S_i\cap S_j\ne\varnothing,\quad
 S_i\nsubseteq S_j,\quad S_j\nsubseteq S_i
 &\quad\text{有交集，双方均有独有实例}.
\end{aligned}
\]
教材模型中，第一类通常按特别法条优先，并先查明有无反向优先等明文处理条款。第三类中，只有本案 \(u\in S_i\cap S_j\) 且可评价为同一个行为，才进入想象竞合的择一重罪讨论；一次同现只能证明交集非空，不能证明包含。特别关系应从要件比对推出，不能把两罪都侵害某一法益当作证明。

\paragraph{例题推导}
甲取得他人唯一的急救药并据为己有，明知断药可能使对方死亡且放任；他人实际因这一风险死亡。题设明确达到盗窃的入罪条件，且不存在独立介入原因。
\[
\begin{aligned}
 &\text{非法占有目的、违背意志取走、达到入罪条件}
       \Longrightarrow E_{\text{盗窃}}(u),\\
 &\text{放任死亡、死亡、因果与结果归责}
       \Longrightarrow E_{\text{故意杀人}}(u).
\end{aligned}
\]
偷走自行车而无人身风险，给出 \(S_{\text{盗窃}}\setminus S_{\text{故意杀人}}\) 的例子；以杀人故意投毒但未取财，给出反方向的例子。本题同一取药单位落在交集且双方无包含，依教材想象竞合模型评价，并按本案适用刑档择重。若甲只预见财产损失，不能从“有死亡”推出杀人故意，第二条推导的主观前提缺失，应另查是否有过失。

\paragraph{基本构成与加重构成}
令 \(B\) 为基本构成，\(A\) 为同一法条明定的升格事实。二者均既遂的典型结构是 \(E_{\text{加}}=B\land A\)，于是 \(E_{\text{加}}\Rightarrow B\)。如第263条的入户抢劫，要分别认定抢劫和入户抢劫的限定事实，才能适用升格刑档。结果加重要再检验基本行为、加重结果的因果与归责以及相应的罪过。形态不一致时应比较完整方案，例如故意杀人中止与故意伤害重伤既遂分别带入相应规则。教材在部分情形采用想象竞合方案，这是形态与评价的解释立场；集合包含只针对已固定的要件与形态，不能自动决定该争议。

\section{第二节\quad 两个行为}
\paragraph{先检索法条，再检验评价是否充分}
对 \(u_1,u_2\) 先独立生成候选 \(E_i(u_1),E_j(u_2)\)，再列出被侵害的对象、法益和结果。令 \(H(u)\) 是这些需评价事实的集合，\(V(r)\) 是方案 \(r\) 在其构成与刑档中已经评价的事实集合。\(H(u_1)\cup H(u_2)\subseteq V(r)\) 可用来检验有无遗漏，不能单凭这一包含式授权合并。对于重叠的量刑事实，也须检查法条是否允许再次评价。多个行为之间的几种典型规则：
\begin{enumerate}
 \item \textbf{结合犯。}明文把独立行为列入一个罪的升格构成。例：拐卖妇女后奸淫该被拐卖妇女，第240条第一款第三项是优先适用的结合规定。对照第157条第二款，走私中以暴力、威胁抗拒缉私，法律明定走私罪与妨害公务罪并罚，不能以二者有联系而覆盖明文。
 \item \textbf{连续犯。}同种行为要先查法条是否以多次为入罪或升格条件，以及有无累计数额的具体规则。教材以财产利益的可替代性解释部分案件的一罪累计，以生命等人身利益的专属性解释数行为、数被害人分别评价的主张；不同人身法益受害人如何计罪有理论和实务差异，不能从“同一罪名”直接推出一罪，也不能从“多人受害”直接推出并罚。
 \item \textbf{吸收犯与不可罚的事后行为。}教材用是否已充分评价来处理前后阶段。如明知是枪支而盗窃并仅顺势持有同一批枪支，按教材吸收观点评价盗窃枪支；若取物时只知普通财物，发现枪支后另起意思非法持有，后来增加的公共安全风险须独立检验。普通处分同一赃物与利用赃物另骗取买家款项也要分开，后者引入新被害人和新损失。
 \item \textbf{牵连犯。}教材模型同时要求两行为有手段与目的或原因与结果关系，并具备类型上通常性；仅有行为人预先安排的目的联系不足以择一。先查特别法条及司法解释，再论教材主张的择重。盗取普通工具后杀人，工具盗窃已独立满足要件时，不能只凭“为了杀人”适用牵连择一。
\end{enumerate}

\paragraph{例题推导}
甲盗窃贵重仿制品 \(u_1\)，随后冒充真品卖给不知情的乙 \(u_2\)。题设明确两罪的数额和主观要件均满足。第一步，所有人的财产失控推出盗窃既遂。第二步，虚构真品、乙错误、付款处分及新增财产损失推出对乙的诈骗既遂。第三步，新增损失 \(h_2\in H(u_2)\) 且 \(h_2\notin V(\text{盗窃评价})\)，排除仅以盗窃吸收；本案无明文合并，按两罪处理。若乙明知是假货，处分错误这一事实为假，不能照搬诈骗结论，仍应另查其他可能的构成。

\section*{记忆与闭卷复述}
记“数行为、列构成、找明文、查新害”。先划 \(U\)，再算每个 \(E_i\)，一行为比集合包含、互斥与交叉，两行为查法定结合或并罚，最后核对新法益与重复评价。用三张卡片分别画包含、交叉和互斥，每次写出两条独有实例或指出为何无法写出。
\begin{enumerate}
 \item \textbf{自测。}甲非法拘禁已经既遂，门锁仍锁住被害人，乙此时故意帮甲守门。能否因已既遂而一律排除共犯？\quad\textbf{核对。}不能。拘禁行为还在继续，查乙对继续拘禁的认识、意思联络与实际参与；“既遂”和“终了”是两条时间线。
 \item \textbf{自测。}举出能推翻“同一目的足以合并两罪”的事实。\quad\textbf{核对。}为杀人独立盗取他人财物作工具，盗窃与杀人分别满足要件。单有共同目的无法给出结合条文，也无法排除独立财产损害。
\end{enumerate}
